\chapter{The Fold: \texttt{iterate\_while} and the Iteration Combinators}
\label{ch:fold}

\chapterorient{This is the keystone of the framework. We met the \emph{fold} once
already, in \cref{ch:lifecycle}, as the shape of the adaptive loop; we promised
then that the same shape would turn out to be the engine of time-stepping, of
Krylov solvers, and of eigenvalue iterations. Now we make good on that promise.
We start from the \texttt{for}-loop every reader knows and recast it, step by
step, as a value carried forward through a single combinator named
\texttt{iterate\_while}. We develop its meaning fully---its \texttt{while}
convention, the one rule about its predicate that everything depends on, its two
sister combinators, and the demand-driven pruning that is the calculus's
signature move. By the end you will be able to read any iterative algorithm in
this book as one specialization of the one combinator taught here.}

\section{From the loop you know to the fold you need}

Every reader of this book has written a \texttt{for}-loop. You set a variable to
a starting value, and then, over and over, you change it: read its current value,
do some work, store a new value back. When a condition fails, you stop and report
what the variable holds. A running sum is the canonical example. In the
imperative style you grew up with, summing a list looks like this:
\begin{lstlisting}[style=l4style]
total = 0
for x in xs:
    total = total + x      -- mutate `total` in place, each pass
return total
\end{lstlisting}
The variable \code{total} is \emph{mutated}: the same storage cell is overwritten
on every pass. \Cref{ch:bigidea} argued at length that mutation is exactly the
property we must give up to run a computation on an immutable-tensor backend. So
we are obliged to find another way to say ``the same thing'' without the
overwrite---and the way we find will turn out to be more powerful than the loop
it replaces.

The trick is to stop thinking of \code{total} as a cell that changes and start
thinking of it as a \emph{value carried forward}. Each pass does not modify a
variable; it \emph{computes a new value from the old one}. ``Carry a value;
take a step that produces the next value; repeat.'' The carried value here is the
running total; the step is ``add the next element''; the repetition walks the
list. Nothing is overwritten---each step's output is a fresh value, and the old
one is simply not referred to again.

This reframing has a name in functional programming: it is a \emph{fold}. A fold
takes a starting value (the \emph{seed}), a step that combines the seed-so-far
with the next thing, and a structure to walk; it threads the accumulated value
through the structure and returns the final accumulation. The running sum is
\code{foldl (+) 0 xs}---``fold the list \code{xs} from the left, starting at
\code{0}, combining with \code{+}.'' \Cref{fig:loopfold} sets the imperative loop
and the fold side by side: the same computation, the mutation traded for a carry.

\begin{figure}[t]
\centering
\begin{tikzpicture}[font=\footnotesize]
  % --- left: the imperative loop, a cell overwritten ---
  \begin{scope}[xshift=0cm]
    \node[font=\small\bfseries, text=simRust] at (0,2.5) {imperative: one cell, overwritten};
    \node[data, minimum width=14mm] (cell) at (0,1.4) {\code{total}};
    \draw[backflow] (cell.east) .. controls (2.2,1.4) and (2.2,2.2) .. (cell.north)
      node[pos=0.5, right=10mm, font=\scriptsize, text=simRust] {$\gets$ overwrite};
    \node[opbox] (add) at (0,0.0) {\code{total = total + x}};
    \draw[flow] (cell.south) -- (add.north);
    \draw[flow] (add.south) -- ++(0,-0.5)
      node[below, font=\scriptsize] {next pass writes \code{total} again};
  \end{scope}
  % divider
  \draw[simGray!50, thick] (4.0,-0.9) -- (4.0,2.7);
  % --- right: the fold, a value threaded forward ---
  \begin{scope}[xshift=8cm]
    \node[font=\small\bfseries, text=simBlue] at (0,2.5) {fold: a value carried forward};
    \node[data] (t0) at (-2.4,1.0) {$0$};
    \node[opbox] (s1) at (-1.0,1.0) {step};
    \node[data] (t1) at (0.2,1.0) {$a_1$};
    \node[opbox] (s2) at (1.4,1.0) {step};
    \node[data] (t2) at (2.6,1.0) {$a_2$};
    \draw[flow] (t0)--(s1); \draw[flow] (s1)--(t1);
    \draw[flow] (t1)--(s2); \draw[flow] (s2)--(t2);
    \draw[flow, simGray] (t2) -- ++(0.9,0) node[right, font=\scriptsize] {$\cdots$};
    \node[font=\scriptsize, text=simGreen] at (-1.0,0.3) {$x_1$};
    \node[font=\scriptsize, text=simGreen] at (1.4,0.3) {$x_2$};
    \draw[flow, simGreen, shorten >=1pt] (-1.0,0.45)--(s1.south);
    \draw[flow, simGreen, shorten >=1pt] (1.4,0.45)--(s2.south);
  \end{scope}
\end{tikzpicture}
\caption[The loop-to-fold bridge]{The same running sum, two ways. On the left,
the imperative loop keeps one storage cell \code{total} and \emph{overwrites} it
each pass---the dashed back-edge is the mutation we must give up. On the right,
the fold threads a \emph{value} forward: each step consumes the carried value
($0$, then $a_1$, then $a_2,\dots$) and the next input $x_i$ and produces a fresh
value. Nothing is overwritten; each arrow carries a distinct immutable value.
This carry-threading picture is the whole chapter in miniature.}
\label{fig:loopfold}
\end{figure}

\subsection{Three folds you already know}

Before we build the simulator's combinator, it pays to fix three patterns from
the functional vocabulary of \cref{ch:language}. Each is a way of walking a
structure; the differences between them are exactly the differences that will
matter for iteration.

\begin{description}[style=nextline, leftmargin=0pt, font=\normalfont\itshape]
  \item[\code{map}] applies a function to each element \emph{independently} and
  collects the results: \lstinline[style=l4style]|map f [x1,x2,x3] = [f x1, f x2, f x3]|.
  The crucial word is \emph{independently}: \code{f x2} does not depend on
  \code{f x1}, so the three applications could be done in any order, or all at
  once. A \code{map} has no carry.
  \item[\code{foldl}] threads an accumulator \emph{left to right}, each step
  depending on the one before:
  \lstinline[style=l4style]|foldl (+) 0 [x1,x2,x3] = ((0 + x1) + x2) + x3|.
  Here order \emph{does} matter: the accumulator entering step~2 is exactly what
  step~1 produced. This dependency is the essence of a fold.
  \item[\code{foldr}] threads from the \emph{right}:
  \lstinline[style=l4style]|foldr (+) 0 [x1,x2,x3] = x1 + (x2 + (x3 + 0))|. For an
  associative operator like \code{+} the two folds agree, but in general they
  bracket the work differently. We will rarely need \code{foldr}; we note it for
  completeness and because it makes the directionality of \code{foldl} vivid.
\end{description}

The distinction between \code{map} and \code{foldl} is the one to carry forward.
A \code{map} is a family of \emph{independent} computations; a \code{foldl} is a
\emph{chain} in which each link consumes the previous. We met exactly this
distinction in \cref{ch:lifecycle}: the electrostatic analysis's several solves
against independent right-hand sides are a \code{map}, while the adaptive loop,
in which each refined mesh comes from the last, is a fold. \Cref{ch:situating}
named it one of the four ``solve corners.'' Iteration---running a step until it
converges---is a fold whose structure is not a fixed list but ``keep going until
a condition holds.''

\begin{keyidea}
A \emph{fold} carries a value forward through a sequence of steps, each step
consuming the value the previous step produced. It is the structure-with-a-carry,
as opposed to \code{map}, the structure-without. Mutation (``overwrite a cell'')
becomes value-threading (``compute the next value''). Every iterative algorithm
in this book is a fold; the rest of the chapter builds the one combinator that
expresses all of them.
\end{keyidea}

\subsection{Folds that compute, not just accumulate}

Summing a list is a fold over a structure that is already laid out---the list
tells us how many steps to take. The folds a simulator runs are different in one
way: \emph{the number of steps is not known in advance}. Newton's method for a
square root, fixed-point iteration, a Krylov solve, a time march---each takes a
step and then \emph{asks whether to take another}, based on the value it now
holds. The structure being folded is not a list of inputs; it is the open-ended
``again?'' of a convergence test.

Two familiar examples make the shape concrete, and we will return to both as
worked examples below.

\emph{Newton's method.} To find $\sqrt{2}$, solve $x^2 - 2 = 0$ by Newton
iteration: from a guess $x$, the next guess is
\[
  x' \;=\; x - \frac{x^2 - 2}{2x} \;=\; \tfrac{1}{2}\!\left(x + \tfrac{2}{x}\right).
\]
We carry $x$, take that step, and repeat until $x$ stops changing (until the
\emph{residual} $|x^2 - 2|$ is small). Carry, step, test, repeat: a fold.

\emph{Fixed-point iteration.} To solve $x = g(x)$ for a contraction $g$, simply
iterate $x' = g(x)$ from any starting point; under a contraction the iterates
converge to the unique fixed point. Again: carry $x$, step by applying $g$, stop
when $x$ barely moves. A fold.

Both share the lifecycle's shape from \cref{sec:fold}: a starting value, a step
that carries state forward, a stopping condition over the accumulated state, a
result threaded through. What they lack---and what a simulator needs---is a
\emph{single named operator} that captures this shape once so that every
algorithm can be written as an instance of it. That operator is
\code{iterate\_while}.

\section{\texorpdfstring{\texttt{iterate\_while}}{iterate\_while}: the canonical primitive}
\label{sec:iteratewhile}

Here is the combinator the whole book turns on. We give its signature first, then
read it slowly, then draw it.

\begin{lstlisting}[style=l4style]
iterate_while
  :: a                                     -- the initial carry
  -> (a -> Bool)                           -- the predicate: keep going?
  -> (a -> { state: a, ...extras })        -- the step: next carry + per-step extras
  -> { final_state: a, trajectory: [{ ...extras }] }
\end{lstlisting}

Read it argument by argument. The first argument is an \emph{initial carry}, a
value of some type \code{a}---it might be a number, a tensor, or a whole record
of solver state. The second is a \emph{predicate}, a pure function from the carry
to a Boolean, answering ``should we take another step?'' The third is the
\emph{step}: from the current carry it produces a record holding the
\code{state}---the \emph{next} carry---together with any per-step
\emph{extras} it wishes to report. The result is a record with two fields: the
\code{final\_state} the carry settled on, and a \code{trajectory}, the list of
all the per-step extras records in order.

\begin{notationbox}
The \lstinline[style=l4style]|{ state: a, ...extras }| notation is a record
(\cref{ch:records}) whose \code{state} field has type \code{a} and which is
\emph{open} in its remaining fields: \code{...extras} stands for ``and whatever
other fields this particular step chooses to emit.'' One step might emit
\lstinline[style=l4style]|{ state: x', residual_norm: r }|; another might emit
\lstinline[style=l4style]|{ state: x' }| with no extras at all. The combinator
does not care which: it forwards the \code{state} as the next carry and piles the
rest into the \code{trajectory}.
\end{notationbox}

In words: \emph{from an initial carry, while a predicate holds of the current
carry, take a step that produces the next carry plus per-step extras, collecting
those extras into a trajectory.} That single sentence is the meaning of
\code{iterate\_while}, and \cref{fig:threaded} draws it. The carry flows left to
right along the top; before each step a predicate gate decides whether to
proceed; the step emits a fresh carry forward and drops its extras downward into
the growing trajectory.

\begin{figure}[t]
\centering
\begin{tikzpicture}[font=\footnotesize, node distance=6mm]
  % carry spine
  \node[data, minimum width=10mm] (a0) at (0,0) {$a_0$};
  \node[diamond, draw=simBlue, fill=simBgBlue, aspect=1.8, inner sep=1pt,
        font=\scriptsize, text=simInk] (p0) at (1.8,0) {$p(a_0)$?};
  \node[opbox, minimum width=12mm] (f0) at (3.9,0) {step};
  \node[data, minimum width=10mm] (a1) at (5.9,0) {$a_1$};
  \node[diamond, draw=simBlue, fill=simBgBlue, aspect=1.8, inner sep=1pt,
        font=\scriptsize, text=simInk] (p1) at (7.7,0) {$p(a_1)$?};
  \node[opbox, minimum width=12mm] (f1) at (9.8,0) {step};
  \node[font=\scriptsize, simGray] (more) at (11.6,0) {$\cdots$};
  % forward carry flow
  \draw[flow] (a0)--(p0);
  \draw[flow] (p0)-- node[above,font=\scriptsize]{yes} (f0);
  \draw[flow] (f0)-- node[above,font=\scriptsize]{\code{.state}} (a1);
  \draw[flow] (a1)--(p1);
  \draw[flow] (p1)-- node[above,font=\scriptsize]{yes} (f1);
  \draw[flow] (f1)-- node[above,font=\scriptsize]{\code{.state}} (more);
  % `no` exits to final_state
  \node[product] (fin) at (7.7,-2.4) {\code{final\_state}};
  \draw[backflow] (p0.south) |- (fin.west);
  \draw[backflow] (p1.south) -- (fin.north);
  % extras drop into trajectory
  \node[data, fill=simBgGold, draw=simGold, minimum width=44mm] (traj) at (6.85,-1.3)
    {\code{trajectory}: $[\,\{\dots e_0\},\ \{\dots e_1\},\ \dots\,]$};
  \draw[refedge] (f0.south) -- (f0.south|-traj.north);
  \draw[refedge] (f1.south) -- (f1.south|-traj.north);
  \node[font=\scriptsize, text=simGold!80!black] at (3.9,-0.85) {extras $e_0$};
  \node[font=\scriptsize, text=simGold!80!black] at (9.8,-0.85) {extras $e_1$};
\end{tikzpicture}
\caption[The threaded carry of \texttt{iterate\_while}]{\texttt{iterate\_while}
as a threaded carry. The carry $a_0, a_1, \dots$ flows forward along the top.
Before each step the predicate gate $p(\cdot)$ tests the \emph{current carry};
on ``yes'' the step runs, emitting the next carry through its \code{.state}
field and dropping its per-step \emph{extras} downward into the growing
\code{trajectory}. The first ``no'' exits the loop, surfacing the current carry
as \code{final\_state}. The predicate fires \emph{first}, before any step---this
is a \texttt{while} loop, not a \texttt{do}/\texttt{while}.}
\label{fig:threaded}
\end{figure}

\subsection{The small-step rule}

We can make the meaning of \code{iterate\_while} exact in one recursive rule,
taken verbatim from the calculus's semantics (\cref{ch:language}). Given a carry
$a$, a predicate $p$, and a step $f$:
\[
\begin{aligned}
\textsf{iterate\_while}\ a\ p\ f \;\to\;\ &\textsf{if}\ p(a)\\
  &\quad\textsf{then}\ \textsf{let}\ \{\textsf{state}\!:\!a',\ \dots e\} = f(a)\ \textsf{in}\\
  &\qquad \textsf{let}\ \{\textsf{final\_state},\ \textsf{trajectory}\} = \textsf{iterate\_while}\ a'\ p\ f\ \textsf{in}\\
  &\qquad \{\textsf{final\_state},\ \textsf{trajectory}\!:\ [\{\dots e\}] \mathbin{+\!\!+} \textsf{trajectory}\}\\
  &\quad\textsf{else}\ \{\textsf{final\_state}\!:\ a,\ \textsf{trajectory}\!:\ [\,]\,\}
\end{aligned}
\]
Read it in prose. Test $p(a)$. If \emph{false}, stop: the answer is $a$ as
\code{final\_state} with an empty trajectory. If \emph{true}, run one step
$f(a)$, obtaining the next carry $a'$ and this step's extras $e$; recurse on
$a'$; and prepend $e$ to whatever trajectory the recursion returns. The recursion
sits in tail position, so an implementation realizes it as an ordinary loop, not
a stack of deferred calls. This is the entire operational content of the
combinator. Everything else in the chapter is a consequence of, or a discipline
around, this one rule.

\subsection{Unrolling the recursion}

The recursive rule is exact, but its shape is easiest to see when we
\emph{unroll} it---write out what it does for a concrete run of three steps
before the predicate finally fails. \Cref{fig:unroll} does exactly that. The
carry $a_0$ passes the predicate and is fed to the first step, producing $a_1$
and extras $e_1$; $a_1$ passes and yields $a_2, e_2$; $a_2$ passes and yields
$a_3, e_3$; and then the predicate fails on $a_3$, so the loop stops. The carry
has walked $a_0 \to a_1 \to a_2 \to a_3$, each link a single step, and the three
extras records have lined up into the trajectory
$[\,e_1, e_2, e_3\,]$ in iteration order.

\begin{figure}[t]
\centering
\begin{tikzpicture}[font=\footnotesize, node distance=8mm]
  % carry chain
  \foreach \i/\lbl in {0/{a_0},1/{a_1},2/{a_2},3/{a_3}}{
    \node[data, minimum width=8mm] (c\i) at (\i*3.0,0) {$\lbl$};
  }
  \foreach \i/\j in {0/1,1/2,2/3}{
    \node[opbox, minimum width=13mm] (st\j) at (\i*3.0+1.5,0) {step};
    \draw[flow] (c\i)--(st\j);
    \draw[flow] (st\j)-- node[above,font=\scriptsize]{\code{.state}} (c\j);
  }
  % predicate annotations above the carries that pass
  \foreach \i in {0,1,2}{
    \node[font=\scriptsize, text=simGreen] at (\i*3.0,0.78) {$p$: yes};
  }
  \node[font=\scriptsize, text=simRust] at (9.0,0.78) {$p$: no\,(stop)};
  % extras drop down to trajectory
  \node[data, fill=simBgGold, draw=simGold, minimum width=58mm] (traj) at (4.5,-1.7)
    {\code{trajectory} $=[\,e_1,\ e_2,\ e_3\,]$};
  \foreach \j/\e in {1/{e_1},2/{e_2},3/{e_3}}{
    \draw[refedge, simGold!80!black] (st\j.south) -- (st\j.south|-traj.north);
    \node[font=\scriptsize, text=simGold!80!black] at ($(st\j.south)+(0,-0.45)$) {$\e$};
  }
  % final_state
  \node[product, minimum width=18mm] (fs) at (9.0,-1.7) {\code{final\_state}\,$=a_3$};
  \draw[flow] (c3.south) -- (fs.north);
\end{tikzpicture}
\caption[A fold unrolled over three steps]{The recursion of
\cref{sec:iteratewhile} unrolled for a run of three steps. The carry threads
$a_0 \to a_1 \to a_2 \to a_3$, each arrow one step. The predicate passes
(\textcolor{simGreen}{yes}) on $a_0, a_1, a_2$ and fails
(\textcolor{simRust}{no}) on $a_3$, ending the loop. Each step drops its
extras downward; the three records line up, in order, as the
\code{trajectory} $[\,e_1, e_2, e_3\,]$, and the final carry $a_3$ becomes
\code{final\_state}. Compare \cref{fig:loopfold}: the carry spine is the same
forward-threaded value, now with extras peeling off below.}
\label{fig:unroll}
\end{figure}

This unrolled picture is the one to hold in mind for the rest of the book.
Whenever a later chapter says ``the solver folds \code{krylov\_step} with
\code{iterate\_while},'' it means precisely \cref{fig:unroll} with ``step''
replaced by the solver's per-step kernel: a carry of solver state threading
forward, a predicate gating each pass, and a residual history peeling off into
the trajectory.

\subsection{What the signature makes structural}
\label{sec:structural}

It is worth pausing to ask why this particular signature, and not some looser
one. An ordinary \texttt{for}-loop lets the body do anything: read any variable,
write any variable, test any condition. The combinator's signature deliberately
forbids most of that, and each thing it forbids buys a property we will rely on.
Four are worth naming.

\begin{description}[style=nextline, leftmargin=0pt, font=\normalfont\itshape]
  \item[The predicate sees the carry only.] Its type
  \lstinline[style=l4style]|a -> Bool| mentions neither the extras nor any
  ambient state. This forces every termination quantity into the carry
  (\cref{sec:predicaterule})---a discipline, but one that makes ``what does this
  loop stop on?'' answerable by reading the carry's type alone.
  \item[The carry is value-threaded; the extras are record-spread.] There is no
  accumulator living \emph{outside} the carry. Anything that must persist across
  iterations is in the carry; anything that is a per-step output is in the extras
  and ends up in the trajectory. The two roles cannot be confused because they
  occupy different slots of the step's return record.
  \item[Per-step output is a declared field, not a side effect.] A step does not
  ``log'' a residual or ``print'' a metric; it \emph{returns} one, as a field of
  its output record. This is what makes the demand-driven pruning of
  \cref{sec:pruning} possible---an output that is a value can be observed or not
  observed, and pruned when not.
  \item[The trajectory is observed structurally.] The list of extras is an
  ordinary value the consumer may read or ignore. Because it is a value and not an
  effect, the calculus can see, at the call site, whether anyone reads it---and
  eliminate the work that produces it when no one does.
\end{description}

None of these is enforced by comments or conventions you must remember; each is a
consequence of the \emph{type}. The signature is doing the work that, in an
imperative loop, the programmer's discipline would have to do---and doing it in a
form a compiler can check and a backend can exploit.

\section{The \texttt{while} convention and the predicate rule}
\label{sec:predicaterule}

Two facts about the predicate are load-bearing, and getting them wrong is the
most common way to misuse the combinator. We state them as rules and then earn
them.

\paragraph{Rule 1: the predicate fires first.} Look again at the small-step
rule. The very first thing it does is test $p(a)$---\emph{before} any step runs.
If the predicate is already false on the initial carry, the body never executes
and the result is the initial carry with an empty trajectory. This is the
\texttt{while} convention (test, then maybe step), not the
\texttt{do}/\texttt{while} convention (step, then test). The distinction is not
pedantry: a solver handed an already-converged input should do zero work, and the
\texttt{while} convention guarantees it.

\paragraph{Rule 2: the predicate is pure on the carry.} The predicate's type is
\lstinline[style=l4style]|a -> Bool|. Its \emph{only} input is the carry. It
cannot read external mutable state, it cannot read the simulator's monadic
environment (\cref{ch:monad}), and---this is the subtle one---it cannot read the
\emph{extras} the step produces. It sees the carry, and nothing else.

\begin{figure}[t]
\centering
\begin{tikzpicture}[font=\footnotesize]
  % --- LEFT: the correct rule ---
  \begin{scope}
    \node[font=\small\bfseries, text=simGreen] at (0,2.4) {correct: predicate reads the carry};
    \node[data, minimum width=22mm, align=center] (carry) at (0,1.3)
      {\code{carry}\\[-1pt]\scriptsize\itshape (incl.\ residual, flag)};
    \node[diamond, draw=simBlue, fill=simBgBlue, aspect=1.6, inner sep=1pt,
          font=\scriptsize] (pred) at (0,-0.2) {\code{p(carry)}};
    \draw[flow] (carry)--(pred);
    \node[font=\scriptsize, text=simGreen] at (2.0,0.55) {one input,};
    \node[font=\scriptsize, text=simGreen] at (2.0,0.2) {no cycle};
  \end{scope}
  \draw[simGray!50, thick] (3.6,-1.2) -- (3.6,2.6);
  % --- RIGHT: the pitfall, crossed out ---
  \begin{scope}[xshift=8.2cm]
    \node[font=\small\bfseries, text=simRust] at (0,2.4) {pitfall: predicate reads the extras};
    \node[opbox, minimum width=16mm] (step) at (-1.3,1.3) {step};
    \node[data, fill=simBgGold, draw=simGold, minimum width=16mm] (ex) at (1.3,1.3) {\code{extras}};
    \node[diamond, draw=simRust, fill=simBgRust, aspect=1.6, inner sep=1pt,
          font=\scriptsize] (badp) at (-1.3,-0.4) {\code{p(extras)}};
    \draw[flow] (step.east)-- node[above,font=\scriptsize]{$e$} (ex.west);
    \draw[flow, simRust] (ex.south) .. controls (1.3,-0.4) and (0.2,-0.4) .. (badp.east);
    \draw[flow, simRust, densely dashed] (badp.north) .. controls (-1.3,0.6) and (-1.3,0.7) .. (step.south)
      node[pos=0.5, left, font=\scriptsize, text=simRust] {gates};
    % the big X
    \draw[simRust, very thick] (-2.4,-1.0) -- (2.4,1.9);
    \draw[simRust, very thick] (-2.4,1.9) -- (2.4,-1.0);
    \node[font=\scriptsize, text=simRust, align=center] at (0,-1.55)
      {circular: the test needs an extra\\the step has not produced yet};
  \end{scope}
\end{tikzpicture}
\caption[The predicate-on-carry rule]{Left: the rule. The predicate reads the
\emph{carry}, which already contains everything termination depends on (a
residual, a converged flag, an iteration count). One input, no cycle. Right: the
pitfall (\cref{sec:predicaterule}). Letting the predicate read the step's
\emph{extras} creates a circular dependency---the gate that decides whether to
run the step needs a value only the step can produce. The fix is to fold any
termination quantity into the carry, so the predicate reads it from there.}
\label{fig:predrule}
\end{figure}

Why must the predicate be pure on the carry? Because termination must be decided
\emph{before} the step runs (Rule~1), and the extras are produced \emph{by} the
step. If the predicate were allowed to read the extras, then on the first
iteration there would be no extras to read---the step that produces them has not
run. That is a circular dependency, drawn crossed-out on the right of
\cref{fig:predrule}. The calculus forbids it structurally: the predicate's type
mentions only \code{a}, so there is simply no way to write a predicate that
reads the extras.

\paragraph{The consequence: fold termination quantities into the carry.} This is
not a limitation; it is a discipline that clarifies every algorithm. Whatever the
loop terminates on---a residual norm below tolerance, a ``converged'' flag, an
iteration count reaching a cap---must live \emph{in the carry}, computed by the
step and read by the predicate on the next pass. The extras hold only the
quantities that are \emph{outputs} (things a monitor or a plot might want); the
carry holds the quantities that are \emph{control inputs} (things the predicate
needs).

A solver that stops when its residual drops below \code{tol} does \emph{not}
write a predicate that inspects ``the residual the step just produced.'' It
carries a \code{residual} field in its state, has the step compute and store it
there, and writes the predicate
\lstinline[style=l4style]|\s -> s.residual > tol|. A solver capped at
\code{max\_it} iterations carries an \code{it} counter in the state and writes
\lstinline[style=l4style]|\s -> s.it < max_it && not s.converged|. The carry is
where convergence lives.

\begin{pitfall}
The single most common mistake with \code{iterate\_while} is trying to let the
predicate test a quantity the step just produced---``loop while the new residual
exceeds the tolerance,'' reaching for the residual as an \emph{extra}. This is a
\emph{circular dependency}: the predicate fires \emph{before} the step
(Rule~1), so on the first pass the extra does not exist yet. The calculus rules
it out: the predicate has type \lstinline[style=l4style]|a -> Bool| and can read
only the carry. \textbf{The fix is always the same}---fold the termination
quantity into the carry. Let the step compute the residual and store it in a
carry field; let the predicate read \emph{that} field. The residual then enters
the next iteration's predicate as part of the carry, exactly when it exists.
\Cref{wex:pitfall} shows the wrong and right versions side by side.
\end{pitfall}

\section{The family: three combinators, one shape}
\label{sec:family}

The combinator of \cref{sec:iteratewhile} is the central member of a small
family. The other two are specializations: one drops a feature it does not need,
one adds a feature for a particular recurrence. \Cref{fig:family} lines all three
up; we take them in turn.

\begin{figure}[t]
\centering
\begin{tikzpicture}[font=\footnotesize, node distance=3mm]
  \def\colw{43mm}
  % --- pure ---
  \node[stage, minimum width=\colw, align=center] (pureH) at (0,0)
    {\code{iterate\_while\_pure}};
  \node[opbox, minimum width=\colw, align=left, below=of pureH] (pureB)
    {carry only,\\no extras\\$\Rightarrow$ returns \code{final\_state}};
  \node[cornertag, below=1mm of pureB] {plain loop};
  % --- base ---
  \node[stage, minimum width=\colw, align=center, right=14mm of pureH] (baseH)
    {\code{iterate\_while}};
  \node[opbox, minimum width=\colw, align=left, below=of baseH] (baseB)
    {carry $+$ extras\\$\Rightarrow$ \code{final\_state}\\$+$ \code{trajectory}};
  \node[cornertag, below=1mm of baseB] {the canonical combinator};
  % --- with_prev ---
  \node[stage, minimum width=\colw, align=center, right=14mm of baseH] (prevH)
    {\code{iterate\_while\_with\_prev}};
  \node[opbox, minimum width=\colw, align=left, below=of prevH] (prevB)
    {carry $+$ extras\\$+$ a threaded \code{prev}\\(two-term recurrence)};
  \node[cornertag, below=1mm of prevB] {bootstrap then steady};
  % arrows showing specialization
  \draw[depedge] (baseH) -- node[above,font=\scriptsize]{$e=()$} (pureH);
  \draw[depedge] (prevH) -- node[above,font=\scriptsize]{\code{prev}$=()$} (baseH);
\end{tikzpicture}
\caption[The iteration-combinator family]{The three members of the iteration
family. The center is \code{iterate\_while}: a carry, a predicate, a step with
extras. Drop the extras ($e=()$) and the trajectory is always empty, leaving
\code{iterate\_while\_pure}, a plain loop returning just the final carry. Add a
second threaded value, the previous iterate, and you get
\code{iterate\_while\_with\_prev}, the driver for two-term recurrences; collapse
its \code{prev} to nothing and it degenerates back to the center. One shape,
three settings.}
\label{fig:family}
\end{figure}

\subsection{\texttt{iterate\_while\_pure}: the plain loop}

Many loops have no per-step output to report---they only thread a carry forward
and want the final value. For these the extras are empty, the trajectory is
always \code{[]}, and carrying it around is noise. The \emph{sugar}
\code{iterate\_while\_pure} drops it:
\begin{lstlisting}[style=l4style]
iterate_while_pure :: a -> (a -> Bool) -> (a -> a) -> a
\end{lstlisting}
The step is now just \lstinline[style=l4style]|a -> a| (next carry, no record),
and the result is just the final carry, no trajectory record to destructure. It
is defined in terms of the general combinator---it \emph{is} the general
combinator with an empty extras record, projected onto its \code{final\_state}:
\begin{lstlisting}[style=l4style]
iterate_while_pure a p f
  = (iterate_while a p (\x -> { state: f x })).final_state
\end{lstlisting}
Newton's method, fixed-point iteration, and a fixed-degree polynomial smoother
all fit \code{iterate\_while\_pure}: each threads a single carry and wants only
the answer. We use it whenever there is nothing per-step worth keeping.

\subsection{\texttt{iterate\_while\_with\_prev}: two-term recurrences}

Some iterations need not only the current carry but the \emph{previous} one. The
conjugate-gradient method (\cref{ch:cg}) builds each search direction from the
current residual \emph{and} a coefficient computed from the previous residual;
a three-term Chebyshev recurrence needs the iterate from two steps back. These
are \emph{two-term recurrences}, and they have an awkward first step: there is no
``previous'' on iteration zero, so the naive code carries a special
\lstinline[style=l4style]|if it == 0| branch inside the step.

The framework removes that branch by splitting the loop into a one-time
\emph{bootstrap} that manufactures the first ``previous'' value, followed by a
\emph{steady} step that is branch-free because it always has a real previous to
work with. The combinator that drives this pattern is
\code{iterate\_while\_with\_prev}:
\begin{lstlisting}[style=l4style]
iterate_while_with_prev
  :: (a -> { state: a, prev: b, ...extras })       -- bootstrap_step (runs once)
  -> a                                             -- initial carry
  -> ((a, b) -> { state: a, prev: b, ...extras })  -- steady_step (carry, prev)
  -> (a -> Bool)                                    -- predicate (carry only)
  -> { final_state: a, trajectory: [{ ...extras }] }
\end{lstlisting}
The \code{bootstrap\_step} fires exactly once, producing both the next carry and
the initial \code{prev}. The \code{steady\_step} then runs the loop, taking the
current carry \emph{and} the threaded \code{prev} as inputs and producing the
next carry, the next \code{prev}, and any extras. The combinator threads
\code{prev} for you; the step never writes the plumbing.

Two disciplines carry over unchanged from the base combinator. The predicate is
\emph{still} pure on the carry alone---it does not read \code{prev}, so any
quantity \code{prev} feeds into termination must be folded into the carry, exactly
as in \cref{sec:predicaterule}. And when no previous value is actually
needed---when \code{prev} can be the empty record---this combinator degenerates
to plain \code{iterate\_while} preceded by one bootstrap step, which is why
\cref{fig:family} draws an arrow collapsing it back to the center. The two are one
family: \code{iterate\_while} is the no-previous idiom, and
\code{iterate\_while\_with\_prev} is its strict generalization for recurrences
that look one step back.

\begin{remark}
Why split the bootstrap out rather than keep the \lstinline[style=l4style]|if it == 0|
test inside the step? Because the split makes the steady step \emph{branch-free}
and the steady carry \emph{one slot lighter}---the previous value lives in the
combinator's threading, not in the state schema. A branch-free step is easier to
reason about, easier to lower to a tensor backend, and free of a conditional that
runs (and is mispredicted) on every iteration but is meaningful on only the
first. \Cref{ch:rotations} studies this ``first-iteration unrolling'' as a
rotation in its own right.
\end{remark}

\section{Demand-driven pruning: the calculus's signature move}
\label{sec:pruning}

We now come to the most elegant property of the combinator, and one of the most
important ideas in the framework. Look again at the step's output record:
\lstinline[style=l4style]|{ state: a, ...extras }|. A step may declare optional
outputs---a residual norm, a breakdown diagnostic, a monitoring metric---\emph{as
record fields, unconditionally}. It always ``produces'' them, in the sense that
they appear in its type. Yet computing a residual norm every iteration is real
work; a production solver that just wants the answer should not pay for
monitoring it will never read.

The naive resolutions to this tension are all familiar and all unpleasant. A
\emph{verbosity flag} threaded through the step (``if \code{verbose}, also compute
the residual''). A \emph{conditional branch} around every optional output. A
\emph{logging monad} (a Writer effect) carrying metrics alongside the state. Each
adds a second code path, and each forces ``monitored'' and ``unmonitored'' to be,
in some sense, different algorithms.

The calculus does none of this. Instead it uses \emph{demand-driven pruning}: the
step declares all its outputs unconditionally, and \emph{whether each output is
actually computed is decided by what a consumer reads}. If a downstream consumer
reads the \code{trajectory}, the per-step extras are computed. If a consumer reads
only \code{final\_state} and never touches the trajectory, then---by the calculus's
evaluation rule---the extras computation is \emph{pruned away entirely}, at every
step, as if it were never written. The same combinator, the same step
definition, the same algorithm: only the consumer's demand differs, and the work
follows the demand.

\begin{keyidea}
Demand-driven pruning makes ``compute residuals for monitoring'' and ``skip
residuals for speed'' the \emph{same algorithm}. The step declares every output
unconditionally; a consumer reads what it wants; outputs nothing reads are
eliminated. There is no verbosity flag, no conditional branch, no logging
monad---there is one step definition, and the demand at the call site decides
what materializes. This is the property that lets the entire book write each
solver \emph{once} and use it for both diagnostics and production.
\end{keyidea}

\subsection{Pruning is graph dead-code elimination}

The mechanism is exactly the dead-code elimination a compiler performs, lifted to
the evaluation of the whole computation graph. The calculus is a
\emph{graph-evaluation language}: every operator produces a record of outputs, and
at solve time only the outputs that some root consumer transitively reaches are
computed. An output no consumer reaches is dead code, and dead code is pruned.

\Cref{fig:pruning} draws this. On the left is the full trajectory: a consumer
reads it, so every step computes its extras and they pile into the list. On the
right, a consumer reads only \code{final\_state}; the trajectory is dead, so the
arrows that would compute and collect the extras are greyed out, and what remains
is the bare carry spine---precisely \code{iterate\_while\_pure}. Formally, when
only \code{final\_state} is observed, the combinator is equivalent to
\begin{lstlisting}[style=l4style]
{ final_state: iterate_while_pure a p f_state, trajectory: [] }
\end{lstlisting}
where \code{f\_state} is the part of the step that computes only the next carry.
The residual-collecting solve and the bare solve are two demands on one
definition.

\begin{figure}[t]
\centering
\begin{tikzpicture}[font=\footnotesize]
  % ============ LEFT: trajectory observed (all live) ============
  \begin{scope}
    \node[font=\small\bfseries, text=simBlue] at (1.5,2.3) {consumer reads \code{trajectory}};
    \foreach \i/\x in {0/0,1/1.5,2/3.0}{
      \node[data, minimum width=7mm] (a\i) at (\x,1.2) {$a_{\i}$};
      \node[opbox, minimum width=7mm, minimum height=5mm] (s\i) at (\x+0.75,1.2) {};
    }
    \node[data, minimum width=7mm] (a3) at (4.5,1.2) {$a_3$};
    \foreach \i in {0,1,2}{ \draw[flow] (a\i.east)--(s\i.west); }
    \draw[flow] (s0.east)--(a1.west); \draw[flow] (s1.east)--(a2.west);
    \draw[flow] (s2.east)--(a3.west);
    % extras live
    \node[data, fill=simBgGold, draw=simGold, minimum width=34mm] (tj) at (2.1,-0.3)
      {\code{trajectory} $=[e_0,e_1,e_2]$};
    \foreach \i in {0,1,2}{ \draw[refedge, simGold!80!black] (s\i.south)--(s\i.south|-tj.north); }
    \node[product, minimum width=18mm] (fs) at (4.5,-0.3) {\code{final\_state}};
    \draw[flow] (a3.south)--(fs.north);
  \end{scope}
  \draw[simGray!50, thick] (5.7,-1.0) -- (5.7,2.5);
  % ============ RIGHT: only final_state observed (extras pruned) ============
  \begin{scope}[xshift=6.6cm]
    \node[font=\small\bfseries, text=simGreen] at (1.5,2.3) {consumer reads only \code{final\_state}};
    \foreach \i/\x in {0/0,1/1.5,2/3.0}{
      \node[data, minimum width=7mm] (b\i) at (\x,1.2) {$a_{\i}$};
      \node[opbox, minimum width=7mm, minimum height=5mm] (t\i) at (\x+0.75,1.2) {};
    }
    \node[data, minimum width=7mm] (b3) at (4.5,1.2) {$a_3$};
    \foreach \i in {0,1,2}{ \draw[flow] (b\i.east)--(t\i.west); }
    \draw[flow] (t0.east)--(b1.west); \draw[flow] (t1.east)--(b2.west);
    \draw[flow] (t2.east)--(b3.west);
    % extras PRUNED (greyed)
    \node[data, fill=white, draw=simGray!40, text=simGray!55, minimum width=34mm] (tjx) at (2.1,-0.3)
      {\textcolor{simGray!55}{\code{trajectory} (pruned)}};
    \foreach \i in {0,1,2}{ \draw[refedge, simGray!35] (t\i.south)--(t\i.south|-tjx.north); }
    \node[product, minimum width=18mm] (fsx) at (4.5,-0.3) {\code{final\_state}};
    \draw[flow] (b3.south)--(fsx.north);
  \end{scope}
\end{tikzpicture}
\caption[Demand-driven pruning]{Demand-driven pruning as graph dead-code
elimination. \emph{Left}: a consumer reads the \code{trajectory}, so every step
computes its extras $e_i$ and they collect into the list---all live. \emph{Right}:
a consumer reads only \code{final\_state}; the trajectory is now dead code, so the
extras computations and the collection arrows (greyed) are eliminated, leaving
only the carry spine. The step definition is identical in both panels; only what
the consumer demands differs. The right panel is exactly
\code{iterate\_while\_pure}.}
\label{fig:pruning}
\end{figure}

This is the calculus's signature move, and it is worth pausing on how much it
buys. In an ordinary imperative solver, the decision ``compute residuals or not''
is a property of the \emph{code}---baked in by a flag or a branch, fixed when the
solver is written. Here it is a property of the \emph{use}: the very same solver,
unmodified, computes residuals exactly when something downstream looks at them.
Monitoring becomes free to add and free to remove. We will lean on this in every
solver chapter; it is why those chapters never present a ``debug version'' beside
a ``fast version.''

\section{What the combinator does, and does not, satisfy}
\label{sec:laws}

A combinator is only as useful as the rewrites it licenses. \code{iterate\_while}
satisfies a handful of laws---facts that let us transform one expression into an
equal one---and, just as importantly, fails to satisfy several tempting ones.
Knowing both keeps us from ``simplifying'' a loop into something that means
something different. We state the most useful affirmatively, and then catalogue
the traps.

\paragraph{The pruning law (the one that earns its keep).} If a consumer observes
only \code{final\_state}, then \code{iterate\_while}~$a~p~f$ equals
\lstinline[style=l4style]|{ final_state: iterate_while_pure a p f_state, trajectory: [] }|,
where \code{f\_state} computes only the step's next carry. This is
\cref{sec:pruning} stated as an equation: the monitored loop and the bare loop are
\emph{equal} expressions whenever nothing reads the trajectory. It is the law
that justifies writing one step definition and using it for both diagnostics and
production.

\paragraph{The base case.} If $p(a)$ is false at the outset,
\code{iterate\_while}~$a~p~f = $
\lstinline[style=l4style]|{ final_state: a, trajectory: [] }| for \emph{any}
step \code{f}. The combinator does not run the body even once before testing---a
direct reading of the small-step rule, and the formal content of Rule~1.

Now the traps. Each is a rewrite that looks innocent and is wrong.

\begin{description}[style=nextline, leftmargin=0pt, font=\normalfont\itshape]
  \item[You may not hoist the predicate out of the loop.] \code{iterate\_while}
  re-evaluates $p$ on \emph{every} carry, not once. Replacing the loop with ``test
  $p(a)$ once, then iterate unconditionally'' turns a terminating loop into an
  infinite one---the predicate would never get the chance to become false. The
  per-iteration re-test is the whole point.
  \item[You may not fuse two loops into one.] Running
  \code{iterate\_while}~$a~p_1~f$ and then \code{iterate\_while} on its result with
  $p_2$ is \emph{not} the same as a single loop with a combined predicate. The two
  phases stop on different conditions and concatenate their trajectories; a single
  flattened loop has one predicate and one trajectory. (This is exactly why GMRES,
  in \cref{ch:gmres}, wraps an \emph{outer} restart loop around its inner
  \code{iterate\_while} rather than flattening the two.)
  \item[You may not swap the predicate and the step.] Testing after the step
  (\texttt{do}/\texttt{while}) is a \emph{different} combinator that can run one
  more time than \code{iterate\_while} on an already-satisfied initial carry. The
  two agree only when the initial predicate is true.
  \item[There is no identity element, and the type does not promise termination.]
  \code{iterate\_while} is a fold, not a monoid operation; there is no carry that
  makes it a no-op for all predicates and steps. And nothing in the signature
  guarantees the loop halts---a step that never lets the predicate go false
  diverges. Termination is an obligation discharged \emph{by the algorithm}: a
  bounded \code{max\_it} folded into the carry, or a convergence proof (Krylov
  methods on well-conditioned systems), or an explicit done-state the step always
  reaches.
\end{description}

The lesson of the non-laws is uniform: \code{iterate\_while} is a faithful record
of an ordered, possibly-unbounded computation, and the order is observable. The
combinator earns its simplicity by \emph{not} pretending iterations commute or
that loops fuse for free---claims that would be convenient and false.

\section{Why one combinator suffices}
\label{sec:onecombinator}

We close the conceptual development by collecting the claim the chapter has been
building toward: \emph{one combinator expresses every iteration in the book}. The
three patterns a simulator runs---independent steps, sequential carries, and
convergence loops---are all \code{iterate\_while} (or its family), seen at
different settings. \Cref{fig:threeuses} shows the three side by side.

\begin{figure}[t]
\centering
\begin{tikzpicture}[font=\footnotesize, node distance=2.5mm]
  \def\cw{38mm}
  % map
  \node[stage, minimum width=\cw] (mH) at (0,0) {as a \textbf{map}};
  \node[opbox, minimum width=\cw, align=left, below=of mH] (mB)
    {predicate counts down\\a fixed list; steps are\\independent (no real carry)};
  \node[cornertag, below=1mm of mB] {independent solves};
  % fold
  \node[stage, minimum width=\cw, right=12mm of mH] (fH) {as a \textbf{fold}};
  \node[opbox, minimum width=\cw, align=left, below=of fH] (fB)
    {carry threads forward;\\each step consumes the\\last (time-stepping)};
  \node[cornertag, below=1mm of fB] {sequential chain};
  % converge
  \node[stage, minimum width=\cw, right=12mm of fH] (cH) {as a \textbf{converge}};
  \node[opbox, minimum width=\cw, align=left, below=of cH] (cB)
    {predicate tests a residual\\folded into the carry;\\stop at tolerance (Krylov)};
  \node[cornertag, below=1mm of cB] {until converged};
  % the unifying brace
  \node[font=\small\bfseries, text=simBlue] (one) at ($(fH)+(0,1.4)$) {one \code{iterate\_while}};
  \draw[depedge] (one) -- (mH.north);
  \draw[depedge] (one) -- (fH.north);
  \draw[depedge] (one) -- (cH.north);
\end{tikzpicture}
\caption[One combinator, three uses]{Three iteration patterns, one combinator.
A \emph{map} of independent steps is \code{iterate\_while} with a predicate that
walks a fixed count and steps that ignore the carry. A \emph{fold} is the
sequential chain---time-stepping, where each step consumes the last. A
\emph{converge} loop tests a residual folded into the carry and stops at a
tolerance---the Krylov solvers. The full catalogue, with the precise sense in
which a map is a degenerate fold, is \cref{ch:combinators}.}
\label{fig:threeuses}
\end{figure}

A \emph{map}---a family of independent solves, like the electrostatic analysis's
several right-hand sides---is \code{iterate\_while} with a predicate that simply
counts off a fixed number of steps and a step whose output does not depend on the
carry. (One ordinarily writes a map as \code{map}, not as a degenerate
\code{iterate\_while}; the point is that it \emph{can} be seen this way, and
\cref{ch:combinators} makes the relationship precise.) A \emph{fold}---a
sequential chain like time-stepping, where each state comes from the last---is
\code{iterate\_while} in its plainest reading. A \emph{convergence loop}---a
Krylov solver running until its residual drops below tolerance---is
\code{iterate\_while} with the residual folded into the carry and the predicate
testing it. Maps, folds, and convergence loops: three faces of one combinator.

\begin{keyidea}
There is \emph{one} iteration combinator. Independent steps (a map), sequential
carries (a fold), and convergence loops (run-until-converged) are all
\code{iterate\_while} at different settings of its predicate and step. The
solver's state-threading, the time-stepper's march, the adaptive loop of
\cref{ch:lifecycle}---every iteration the simulator performs is one specialization
of the combinator built in this chapter. Learn it once and you have learned the
control structure of the entire simulator.
\end{keyidea}

The headline instance is the Krylov step. \Cref{ch:cg} and \cref{ch:gmres} build
the conjugate-gradient and GMRES solvers, and at their heart each is a single
\code{krylov\_step} folded by \code{iterate\_while}: the carry is the solver's
working state (current iterate, residual, search direction, convergence flag),
the step applies the operator once and updates that state, the predicate reads
the convergence flag and iteration count from the carry, and the trajectory---when
a consumer asks for it---is the residual history. Everything those chapters do is
filling in the step; the loop is the one we have here. \Cref{ch:combinators}
catalogues the full family of specializations, and \cref{ch:monad} explains how
the step threads the simulator's monadic state while the carry stays a pure value.

\section{Worked examples}

\begin{workedexample}{Newton's method for $\sqrt{2}$ as a fold}{newton}
We compute $\sqrt 2$ by Newton iteration, written as \code{iterate\_while}. The
carry holds both the current estimate and the residual that the predicate will
test---folding the termination quantity into the carry, exactly as
\cref{sec:predicaterule} requires.

The carry is a record \lstinline[style=l4style]|{ x: Float, residual: Float }|.
The step applies the Newton update and recomputes the residual; the predicate
tests the residual against a tolerance.
\begin{lstlisting}[style=l4style]
newton_sqrt2 :: Float -> { final_state: { x, residual }, trajectory: [...] }
newton_sqrt2 x0 =
  iterate_while
    { x: x0, residual: abs (x0*x0 - 2) }        -- initial carry
    (\s -> s.residual > 1e-12)                  -- predicate: pure on the carry
    (\s -> let x' = 0.5 * (s.x + 2 / s.x)       -- Newton step
           in  { state:    { x: x', residual: abs (x'*x' - 2) },
                 residual: abs (x'*x' - 2) })   -- extra: the per-step residual
\end{lstlisting}
Starting from $x_0 = 1$, trace the carry. Each step computes
$x' = \tfrac12(x + 2/x)$ and the residual $|x'^2 - 2|$:
\begin{center}
\setlength{\tabcolsep}{8pt}
\renewcommand{\arraystretch}{1.15}
\begin{tabular}{@{}clll@{}}
\toprule
step & carry $x$ & next $x' = \tfrac12(x + 2/x)$ & residual $|x'^2 - 2|$ \\
\midrule
0 & $1.000000000000$ & $1.500000000000$ & $2.5\times10^{-1}$ \\
1 & $1.500000000000$ & $1.416666666667$ & $6.9\times10^{-3}$ \\
2 & $1.416666666667$ & $1.414215686275$ & $6.0\times10^{-6}$ \\
3 & $1.414215686275$ & $1.414213562374$ & $4.5\times10^{-12}$ \\
4 & $1.414213562374$ & $1.414213562373$ & $<10^{-12}$ \;(stop) \\
\bottomrule
\end{tabular}
\end{center}
The residual falls roughly quadratically---each step doubles the correct digits.
After step~4 the predicate \lstinline[style=l4style]|s.residual > 1e-12| is false,
the loop halts, and \code{final\_state.x} is $\sqrt 2$ to twelve places. Two
points to notice. First, the predicate never touches the freshly computed
residual directly; it reads \code{s.residual} from the carry, which the
\emph{previous} step stored there. Second, the residual appears \emph{both} in the
carry (so the predicate can read it) and in the extras (so a consumer plotting
convergence can read the trajectory)---two readers, one computed value.
\end{workedexample}

\begin{workedexample}{A geometric-series partial sum as \texttt{iterate\_while\_pure}}{geom}
When a loop has nothing per-step to report, \code{iterate\_while\_pure} is the
idiom. We sum the geometric series $\sum_{k=0}^{\infty} r^k = 1/(1-r)$ for
$|r| < 1$ by accumulating partial sums until a new term changes the sum
negligibly.

The carry is a record \lstinline[style=l4style]|{ sum: Float, term: Float }|: the
partial sum so far and the current term $r^k$. The step adds the term and
advances it to $r^{k+1}$; the predicate stops when the term is below tolerance.
There are no extras, so we use the pure form and read back only the final sum.
\begin{lstlisting}[style=l4style]
geom_sum :: Float -> Float
geom_sum r =
  (iterate_while_pure
     { sum: 0.0, term: 1.0 }                    -- start: sum 0, term r^0 = 1
     (\s -> abs s.term > 1e-10)                 -- stop when the term is tiny
     (\s -> { sum: s.sum + s.term, term: s.term * r })   -- accumulate, advance
  ).sum
\end{lstlisting}
For $r = \tfrac12$ the exact sum is $1/(1-\tfrac12) = 2$. Tracing the carry:
\begin{center}
\setlength{\tabcolsep}{8pt}
\renewcommand{\arraystretch}{1.1}
\begin{tabular}{@{}clll@{}}
\toprule
step & \code{term} ($= r^k$) & \code{sum} after adding & remaining error $2 - \text{sum}$ \\
\midrule
0 & $1.0000$    & $1.0000$ & $1.0000$ \\
1 & $0.5000$    & $1.5000$ & $0.5000$ \\
2 & $0.2500$    & $1.7500$ & $0.2500$ \\
3 & $0.1250$    & $1.8750$ & $0.1250$ \\
$\vdots$ & $\vdots$ & $\vdots$ & $\vdots$ \\
33 & $\approx10^{-10}$ & $\approx 2.0000$ & $<10^{-10}$ \;(stop) \\
\bottomrule
\end{tabular}
\end{center}
The carry threads two values, the predicate reads one of them (\code{term}), and
because nothing per-step is collected the \code{iterate\_while\_pure} sugar saves
us a trajectory we would never read. This is fixed-point iteration in disguise:
the map $s \mapsto (\,s.\code{sum} + s.\code{term},\; r\cdot s.\code{term}\,)$ is a
contraction, and the loop walks to its fixed point.
\end{workedexample}

\begin{workedexample}{The predicate-on-extras pitfall, wrong and right}{pitfall}
This example shows the mistake of \cref{sec:predicaterule} concretely and its fix
side by side. We want a generic ``iterate until the residual is below
\code{tol}'' loop. The tempting but \emph{wrong} sketch tries to let the predicate
read the residual the step produces as an extra:
\begin{lstlisting}[style=l4style]
-- WRONG: the predicate wants to read an extra the step has not produced yet.
solve_BAD x0 =
  iterate_while_BAD
    x0
    (\extras -> extras.residual > tol)       -- (!) reads the step's extra
    (\x -> let x'  = step_body x
               r   = residual_of x'
           in  { state: x', residual: r })   -- residual is an *extra*
\end{lstlisting}
The predicate \lstinline[style=l4style]|\extras -> extras.residual > tol| needs an
\code{extras} record to inspect---but the predicate fires \emph{before} the step
(Rule~1), so on the first iteration no such record exists. The dependency is
circular; \cref{fig:predrule} draws it crossed out. The calculus will not even
let us write this: the predicate's type is \lstinline[style=l4style]|a -> Bool|,
so it cannot name an \code{extras} argument at all.

The fix folds the residual into the carry. The carry becomes a record carrying
both the iterate and its residual; the step computes the residual and stores it
in the carry; the predicate reads it from there:
\begin{lstlisting}[style=l4style]
-- RIGHT: the residual lives in the carry; the predicate reads it from there.
solve_OK x0 =
  iterate_while
    { x: x0, residual: residual_of x0 }      -- residual folded into the carry
    (\s -> s.residual > tol)                 -- predicate: pure on the carry
    (\s -> let x' = step_body s.x
               r  = residual_of x'
           in  { state: { x: x', residual: r },   -- carry the residual forward
                 residual: r })                   -- also surface it as an extra
\end{lstlisting}
Now the predicate reads \code{s.residual}, a field of the carry that the previous
step stored---which exists from the start because the \emph{initial} carry
computes it too. The residual is still available to trajectory consumers as an
extra; it is simply \emph{also} in the carry so the predicate can see it. The
single change---move the control quantity from the extras into the carry---turns
the impossible loop into the canonical one. This is the move to reach for every
time a predicate seems to need something the step produced.
\end{workedexample}

\begin{example}[A conjugate-gradient step, previewed]
We can already sketch the shape \cref{ch:cg} fills in. A conjugate-gradient solve
is \code{iterate\_while} whose carry is the CG working state---current iterate
$\vx$, residual $\vr$, search direction $\vp$, residual norm, and a
\code{converged} flag---and whose step is one \code{krylov\_step}: apply the
operator once, update $\vx$, $\vr$, $\vp$, recompute the residual norm, and set
\code{converged} when it drops below tolerance. The predicate reads the carry:
\lstinline[style=l4style]|\s -> s.it < max_it && not s.converged|. The trajectory,
when a consumer reads it, is the residual history. The entire control structure is
the combinator of this chapter; \cref{ch:cg} supplies only the step body.
\end{example}

\summarysection
An imperative loop overwrites a cell; a \emph{fold} carries a value forward,
each step computing the next from the last. This single reframing---mutation into
value-threading---turns every iteration in the book into an instance of one
combinator, \code{iterate\_while}. Its signature is
\lstinline[style=l4style]|a -> (a -> Bool) -> (a -> { state: a, ...extras }) -> { final_state, trajectory }|:
from an initial carry, while a predicate holds of the current carry, take a step
that produces the next carry plus per-step extras, collecting the extras into a
trajectory. Two rules about the predicate are load-bearing: it fires
\emph{first} (the \texttt{while} convention), and it is \emph{pure on the carry}
(it sees the carry and nothing else). The consequence is the central discipline:
any termination quantity---a residual, a flag, a count---must be \emph{folded into
the carry}, computed by the step and read by the predicate. Trying instead to test
an \emph{extra} the step produced is the predicate-on-extras pitfall, a circular
dependency the calculus forbids by typing. The family has three members:
\code{iterate\_while\_pure} (no extras, returns just the final carry),
\code{iterate\_while} (the canonical form), and
\code{iterate\_while\_with\_prev} (a bootstrap plus a branch-free steady step for
two-term recurrences). \emph{Demand-driven pruning}---graph dead-code
elimination---makes a monitored solve and a bare solve the \emph{same algorithm}:
the step declares every output unconditionally, and outputs no consumer reads are
eliminated, with no flag, branch, or logging monad. Finally, one combinator
suffices: maps, folds, and convergence loops are all \code{iterate\_while} at
different settings, and the headline instance, the Krylov step
(\cref{ch:cg},~\cref{ch:gmres}), is a single step folded by exactly this loop.

\furtherreading
The fold and its relatives are foundational in functional programming; Peyton
Jones's account of Haskell~\citep{peytonjones2003} develops \code{map},
\code{foldl}, \code{foldr}, and the laws relating them. The idea that
``effects''---here, the optional per-step outputs---can be expressed as values
threaded through pure code, rather than as ambient mutation or a logging channel,
is the subject of Wadler's monads tutorial~\citep{wadler1995monads}; we take up the
state monad that threads the simulator's own mutable bookkeeping in
\cref{ch:monad}. For the numerical iterations the combinator will drive---Krylov
methods, their convergence, and the residual histories the trajectory
records---Saad's monograph~\citep{saad2003} is the standard reference, and we build
those solvers as folds in \cref{ch:cg} and \cref{ch:gmres}.

\exercisessection
\begin{enumerate}[nosep]
  \item Write the running sum of a list as both an imperative \texttt{for}-loop and
  a \code{foldl}. Identify, in each, the carry and the step. Which version could
  you run in parallel, and which not? (Hint: which one has a carry?)
  \item In \cref{wex:newton} the residual appears in two places in
  the step's output record---in the carry's \code{residual} field and as a
  top-level \code{residual} extra. Explain what each copy is for, and which
  reader consumes which. What would break if you kept only the extra?
  \item State the \texttt{while} convention precisely and contrast it with
  \texttt{do}/\texttt{while}. Construct an initial carry on which the two
  conventions give \emph{different} numbers of steps. Why does a solver want the
  \texttt{while} convention?
  \item A colleague proposes a predicate of type
  \lstinline[style=l4style]|(a, extras) -> Bool| so the loop can ``stop as soon
  as the step's residual is small.'' Explain the circular dependency this creates
  on the first iteration, and rewrite their loop using the carry-folding fix of
  \cref{sec:predicaterule}.
  \item Using the small-step rule of \cref{sec:iteratewhile}, evaluate
  \lstinline[style=l4style]|iterate_while a (\_ -> False) f| by hand for any step
  \code{f}. What is the \code{final\_state}? The \code{trajectory}? How many times
  does \code{f} run? Which of the two predicate rules does this exercise probe?
  \item A solver is written with a step that always computes a residual norm as an
  extra. In one run a consumer plots the residual history; in another it reads only
  the final iterate. Using \cref{sec:pruning}, describe what computation each run
  actually performs, and argue that the two are running the \emph{same} algorithm.
  Why is this preferable to a \code{verbose} flag?
  \item Express a fixed-degree polynomial smoother---a loop that runs exactly $d$
  steps regardless of any residual---as \code{iterate\_while\_pure}. What does the
  carry contain? What does the predicate test? (Hint: fold a step counter into the
  carry.)
  \item \emph{(Two-term recurrence.)} The Fibonacci sequence obeys
  $F_{k+1} = F_k + F_{k-1}$, a two-term recurrence. Sketch its computation using
  \code{iterate\_while\_with\_prev}: identify the bootstrap step (what is the first
  ``previous''?), the steady step, the carry, and the predicate. Then argue that
  the same sequence can be computed by plain \code{iterate\_while} if you widen the
  carry to hold \emph{both} $F_k$ and $F_{k-1}$. What does
  \code{iterate\_while\_with\_prev} buy over that widening?
  \item \emph{(Synthesis.)} The adaptive loop of \cref{ch:lifecycle} threads the
  mesh forward, refining it each pass until the error is small enough. Cast it as
  \code{iterate\_while}: name the carry, the step, the predicate, and the
  termination quantity that must be folded into the carry. Is its trajectory ever
  read? If not, what does \cref{sec:pruning} say happens to the per-step extras?
\end{enumerate}
